%%%%%
%
%% Dans ce fichier, les �ventuels accents sont cod�s
%% en ISO-8859-15 (aka latin9)
%% La ligne suivante donne trois exemples 
%% e accent aigu : �, e accent grave : �, e-dans-l'o : �
\section{�chantillonnage et moyenne harmonique d'une variable sur une nouvelle base
de temps donn�e par un time\_{bnds} en tenant compte de valeurs invalides}

\itemstyle{Op\'erateur~:} \operfmt{MoyHarmInvEch\_tps\_bnds}
\index{MoyHarmInvEch\_tps\_bnds}

%%%%%OBJET%%%%
\itemstyle{Objet~:}
Cette fonction calcule, pour chaque instant $t$ d'une
base de temps $time_{out}$, 
une variable $X_{out}$, contenant la moyenne harmonique des
valeurs de  $X_{in}$ entre les bornes donn�es par le $time_{bnds}$ associ� � $time_{out}$.
Le calcul n'est r�alis� que pour des valeurs < $val_{inv}$.
Aucune supposition n'est faite sur la r�gularit� de
l'�chantillonnage des 2 variables. Il est cependant n�cessaire
que $time_{in}$ (le temps de $X_{in}$) soit "� peu pr�s" tri� par ordre
croissant.


%%%%%ARGUMENTS%%%%
\bigskip
%% Dans ce fichier, les �ventuels accents sont cod�s
%% en ISO-8859-15 (aka latin9)
%% La ligne suivante donne trois exemples 
%% e accent aigu : �, e accent grave : �, e-dans-l'o : �
%%%%%ENTREES%%%%
\begin{tabarg}
\multicolumn{3}{l}{\titretab {Entr\'ees}}\\
\hline
$X_{in}$ & Vecteur, \texttt{float} & Variable d'origine\\
$toler_{mnq}$ & Coeff., \texttt{Int\_32} & Seuil de calcul de la moyenne harmonique(nb. �chant). A chaque instant de la base de temps d'arriv�e, 
la moyenne harmonique est calcul�e si le nombre d'�chantillons dans l'intervalle de temps est strictement sup�rieur � $toler_{mnq}$. Sinon, 
le r�sultat est "valeur manquante". Si $toler_{mnq}$ = 0, le calcul est fait d�s qu'il y a au moins une valeur valide dans l'�chantillon. \\
$time_{out}$ & Vecteur, \texttt{double} & Base de temps d'arriv�e (secondes)\\
$val_{inv}$ & Coeff., \texttt{float} & valeur invalide pour laquelle la Moyenne harmonique n'est pas calcul�e \\
$val_{remplac}$ & Coeff., \texttt{float} & valeur de remplacement quand la valeur atteint val\_{inv} %avec une tol�rance 
au moins un fois dans l'intervalle\\
\end{tabarg}

\medskip
%%%%%SORTIES%%%%
\begin{tabarg}
\multicolumn{3}{l}{\titretab {Sorties}}\\
\hline
$X_{out}$ & Vecteur, \texttt{float} & $X_{in}$ moyenn� et �chantillonn� sur
la base de $time_{out}$ (m�me unit� que $X_{in})$\\
\end{tabarg}



%%%%%FORMULE%%%%
\itemstyle{Formule}

%% Formule de l'operateur : MoyHarmInvEch_tps_bnds
%% Dans ce fichier, les �ventuels accents sont cod�s
%% en ISO-8859-15 (aka latin9)
%% La ligne suivante donne trois exemples 
%% e accent aigu : �, e accent grave : �, e-dans-l'o : �
Pour  i=$0$, Nbre  $Time_{bnds}$ :
%%%% tolerance_inv n'est pas impl�ment�e dans cet op�rateur
%\[
%Proportion_{valid}=\frac{\texttt{nb de }X_{in}[j] < %val_{inv}}{\texttt{nb de valeur dans l'intervalle } %Time\_bnds}
%\]

\[
X_{out}[i] =
\begin{cases}
\begin{tabular}{ll}
$
\frac{N_{i}}
{\displaystyle{\sum_{tb[i][0]}^{tb[i][1]} \frac{1}{X_{in}[j]}} }
$
&
\texttt{ si tous les terme de l'intervalle sont $<  val_{inv}$}
\\
\texttt{$val_{replac}$} & \texttt{ si $val_{inv}$ est atteint au moins un fois}
\end{tabular} 
\\
\end{cases}
\]

Pour chaque $X_{in}[j]$ inf\'erieur \`a $val_{inv}$ avec $N_i$ le nombre de donn\'ees dans l'intervalle tb  pour lesquels $X_{in}[j]$ est inf\'erieur \`a $val_{inv}$ .


%%%%COMMENTAIRE%%%%
\itemstyle{Commentaire~:}
\textit{tb} est un tableau d'intervalles de temps d�termin�s � partir du vecteur temps $time_{out}$.

%%%%%AUTEUR%%%%
\itemstyle{Auteur~:}
CNRM/GMEI/TRAMM

%%%%%REFERENCE%%%%
\itemstyle{R\'ef\'erences~:}
N�ant.
